% =========================================================
%  config/listings.tex
%  Gaya kode Python dan C#.
% =========================================================

\lstdefinelanguage{CSharp}{
  morekeywords={abstract,as,base,bool,break,byte,case,catch,char,class,const,
    continue,decimal,default,delegate,do,double,else,enum,event,explicit,
    extern,false,finally,fixed,float,for,foreach,goto,if,implicit,in,int,
    interface,internal,is,lock,long,namespace,new,null,object,operator,out,
    override,params,private,protected,public,readonly,ref,return,sbyte,
    sealed,short,sizeof,stackalloc,static,string,struct,switch,this,throw,
    true,try,typeof,uint,ulong,unchecked,unsafe,ushort,using,virtual,void,
    volatile,while,var,get,set},
  sensitive=true,
  morecomment=[l]{//},
  morecomment=[s]{/*}{*/},
  morestring=[b]",
}

\lstdefinestyle{modulkode}{
  basicstyle=\ttfamily\footnotesize\color{ubCodeText},
  keywordstyle=\color{ubCodeKey}\bfseries,
  commentstyle=\color{ubCodeCom}\itshape,
  stringstyle=\color{ubCodeStr},
  numberstyle=\tiny\color{ubCodeNum},
  numbers=left,
  numbersep=8pt,
  backgroundcolor=\color{ubCodeBg},
  frame=none,
  showstringspaces=false,
  breaklines=true,
  breakatwhitespace=false,
  postbreak=\mbox{\textcolor{ubCodeKey}{$\hookrightarrow$}\space},
  tabsize=4,
  columns=fullflexible,
  keepspaces=true,
  xleftmargin=14pt,
  xrightmargin=3pt,
  aboveskip=6pt,
  belowskip=6pt,
  upquote=true,
}
\lstset{style=modulkode}

% \tampilkode[opsi]{berkas}{keterangan}
\newcommand{\tampilkode}[3][]{%
  \Needspace{6cm}%
  \lstinputlisting[style=modulkode, caption={#3}, #1]{#2}%
}

% \tampilkodebaris[bahasa]{berkas}{awal}{akhir}{keterangan}
\newcommand{\tampilkodebaris}[5][Python]{%
  \Needspace{5.5cm}%
  \lstinputlisting[style=modulkode, language=#1,
    firstline=#3, lastline=#4, caption={#5}]{#2}%
}
